\documentclass[letterpaper,10pt,conference]{ieeeconf}
\IEEEoverridecommandlockouts
\overrideIEEEmargins

\usepackage{amsmath,amssymb,bm}
\usepackage{booktabs}
\usepackage{graphicx}
\usepackage{multirow}
\usepackage{url}
\usepackage{cite}
\usepackage{xcolor}
\usepackage{tikz}
\usetikzlibrary{arrows.meta,positioning,fit}

% Citation/link blue sampled from the CoSAM manuscript (#367DBD).
\definecolor{cosamLinkBlue}{HTML}{367DBD}

% ieeeconf defines an obsolete Natbib compatibility hook.  Removing it before
% loading hyperref restores clickable numerical citations.
\makeatletter
\let\NAT@parse\undefined
\makeatother
\usepackage[
  colorlinks=true,
  linkcolor=cosamLinkBlue,
  citecolor=cosamLinkBlue,
  urlcolor=black
]{hyperref}

\definecolor{todoRed}{RGB}{176,36,36}
\definecolor{inputBlue}{RGB}{225,238,250}
\definecolor{geometryGreen}{RGB}{224,243,231}
\definecolor{routeOrange}{RGB}{253,237,214}
\definecolor{outputPurple}{RGB}{238,229,247}
\newcommand{\todo}[1]{\textcolor{todoRed}{\textbf{[TODO: #1]}}}
\newcommand{\tbd}{\textcolor{todoRed}{--}}
\newcommand{\R}{\mathbb{R}}
\newcommand{\I}{\mathbf{I}}
\newcommand{\D}{\mathbf{D}}
\newcommand{\K}{\mathbf{K}}
\newcommand{\T}{\mathbf{T}}
\newcommand{\M}{\mathbf{M}}
\newcommand{\F}{\mathbf{F}}
\newcommand{\U}{\mathbf{U}}
\newcommand{\n}{\mathbf{n}}
\newcommand{\od}{\mathbin{\odot}}
\newsavebox{\roboticplatformbox}

\title{\LARGE \bf
DistriNVS: Distributional Geometry Transport and Visibility-Constrained Completion for Cross-Endoscope Novel-View Synthesis}

\author{Anonymous Authors}

\begin{document}
\maketitle
\thispagestyle{empty}
\pagestyle{empty}

\begin{abstract}
Cross-endoscope novel-view synthesis must both transport tissue observed by
the source camera and infer tissue revealed only from the target viewpoint.
These two regimes are confounded when sparse or erroneous depth is projected
as deterministic points: sampling gaps, depth failures, surface collisions,
and genuine disocclusions all become visually identical holes. We present
DistriNVS, a reliability-aware framework that models every source depth as a
distribution and propagates it through calibrated geometry. Its
distributional surface splatting renderer produces an anisotropic target-plane
footprint together with support, depth variance, and collision entropy. These
statistics drive a transport-or-synthesize router: a residual transport expert
corrects co-visible appearance, whereas an uncertainty-conditioned Fourier
network synthesizes unsupported or geometrically unreliable content. A final hard composition prevents the
network from modifying trustworthy projected pixels. Joint target RGB-depth
prediction, backward reprojection, and an explicit risk head provide geometric
and failure-awareness constraints. The prescribed protocol trains on a fixed
set of 13 iMED scenes and evaluates without adaptation on held-out EndoVis
datasets 8 and 9 using a strict raw-overlap protocol. 
% \todo{Insert the final test-set
% sample count, the strongest quantitative improvements with confidence
% intervals, and runtime after all experiments are frozen.}
\end{abstract}

\section{Introduction}
\label{sec:introduction}

Robot-assisted minimally invasive surgery depends on endoscopic vision for
scene perception and instrument guidance
~\cite{wei2026roboticperception,sunmola2026surgicalgaussiansurfels}.
Cross-endoscope novel-view synthesis can broaden this visual context by
predicting a target-endoscope view from synchronized appearance and geometry
acquired by a spatially displaced source endoscope~\cite{wang2026endoe2egs}.
Unlike left-to-right synthesis within one stereo endoscope, the target may
observe anatomy absent from both source cameras. The task therefore combines
two epistemically distinct regimes: co-visible anatomy has measured appearance
that should be transported through calibrated geometry, whereas target-only
anatomy lacks source evidence and must be synthesized. A reliable system should
preserve supported measurements, identify inferred content, and thereby support
wider-field robotic perception, retrospective visualization, and surgical
simulation~\cite{hsieh2026extrags}.
Such uncertainty-aware cross-view completion may also support visual-state
construction for surgical robot world models~\cite{he2025surgworld}.

Achieving this separation is difficult because endoscopic depth and calibration
are unreliable around specularities, weak texture, depth discontinuities,
instruments, and deforming tissue~\cite{allan2021scared}. Conventional forward
warping represents each valid source pixel by one 3-D point; discretized
projection leaves sampling gaps, while multiple points may compete for the same
target pixel. A binary validity mask then conflates missing depth, rasterization
holes, depth-ordering conflicts, and genuine cross-endoscope disocclusions.
Treating these cases uniformly either degrades observed structure or permits
hallucination where geometric evidence remains trustworthy.

\begin{figure}[t]
\centering
\vspace*{2mm}
\sbox{\roboticplatformbox}{%
  \includegraphics[width=\columnwidth]{figures/distrinvs_teaser/distrinvs_teaser.pdf}%
}
\begin{tikzpicture}
  \clip[rounded corners=0mm]
    (0,0) rectangle
    (\wd\roboticplatformbox,\ht\roboticplatformbox);
  \node[anchor=south west,inner sep=0pt] at (0,0)
    {\usebox{\roboticplatformbox}};
\end{tikzpicture}
\caption{Cross-endoscope synthesis with explicit provenance. DSS transports
source-supported anatomy and exposes visibility uncertainty; DistriNVS
synthesizes target-only content while trusted composition preserves reliable
projected RGB-D.}
\label{fig:conditional_refinement}
\vspace{-0.4cm}
\end{figure}

Existing methods address complementary parts of this problem but do not expose
the required evidence boundary. Differentiable point rendering and softmax
splatting improve feature transport and collision handling
~\cite{wiles2020synsin,niklaus2020softmax}, while reference-guided inpainting
combines depth-based alignment with learned completion~\cite{zhao2023geofill}.
Large-mask inpainting networks provide strong global image priors
~\cite{suvorov2022lama,chu2023uffc}, and endoscopic video completion adds
depth, temporal, or stereo cues
~\cite{daher2023temporal,zhang2024daevi,zou2025ssifnet}. Endoscopic neural
rendering has likewise advanced sparse-view reconstruction and uncertainty
modeling~\cite{li2024endosparse,guo2025ucnerf}. However, deterministic
projection commonly assigns equal reliability to all valid depths, while
image-space uncertainty does not explicitly describe how source geometry
supports a spatially offset target view. Neither approach propagates per-pixel
depth uncertainty through the cross-endoscope transform and uses the resulting
visibility statistics to constrain where appearance may be synthesized.

Our central premise is to \emph{transport reliable geometric evidence and
synthesize only unsupported or geometrically unreliable regions}. We introduce
DistriNVS, a reliability-aware framework
for direct cross-endoscope novel-view synthesis. Its distributional surface
splatting (DSS) represents each source depth as a distribution whose bounded
mean and scale, surface normal, and confidence are predicted by a depth
reliability encoder. DSS projects deterministic sigma samples through the
calibrated cross-endoscope transform and aggregates source features using
Jacobian- and normal-aware footprints with soft nearest-surface visibility,
yielding target support, depth variance, and collision entropy. These statistics
drive a continuous gate between residual transport for co-visible regions and
uncertainty-conditioned Fourier completion for unsupported or ambiguous
regions. A trusted mask hard-copies trusted DSS-supported RGB-D into the output,
while target-depth and risk prediction plus backward RGB-D reprojection constrain
completion during training.

The paper makes three contributions:
\begin{itemize}
    \item We formulate cross-endoscope novel-view synthesis as a
    reliability-aware separation between transport of anatomy observed by the
    source endoscope and constrained inference of content revealed only by the
    target endoscope.
    \item We introduce DSS, a differentiable renderer that propagates depth
    distributions through calibrated cross-endoscope geometry and exposes
    target support, depth variance, and collision entropy beyond binary validity.
    \item We design a two-expert decoder that routes co-visible regions to
    residual transport and unsupported or geometrically unreliable regions to uncertainty-conditioned
    Fourier completion, while trusted-pixel composition preserves reliable
    projected RGB-D.
\end{itemize}

\section{Related Work}
\label{sec:related}

We briefly review image completion, geometry-guided view transport, and
endoscopic neural rendering, with particular attention to how existing
approaches represent missing content, exploit calibrated geometry, and quantify
uncertainty.

\begin{figure*}[t]
\centering
\includegraphics[width=\textwidth]{figures/distrinvs_framework/distrinvs_framework.pdf}
\caption{DistriNVS framework. A reliability encoder predicts source geometry
and features. DSS transports deterministic depth hypotheses through the
calibrated cross-endoscope transform and jointly produces warped evidence and
target-plane visibility statistics. These statistics condition a bounded
transport-or-synthesize gate, while an independent trusted mask hard-copies
accepted DSS evidence. Solid arrows denote inference; the dashed lower lane is
used only during training.}
\label{fig:overview}
\end{figure*}

\subsection{Image and Endoscopic Video Inpainting}

Generic image inpainting combines mask-aware feature extraction with
increasingly broad appearance priors. Partial and gated convolutions prevent
invalid activations from contaminating local features
~\cite{liu2018partial,yu2019gated}, whereas LaMa and UFFC use Fourier operations
for global context and MAT restricts attention to valid tokens
~\cite{suvorov2022lama,chu2023uffc,li2022mat}. Recent diffusion models such as
BrushNet and TurboFill incorporate masked-image features into pretrained
generators or adapt few-step sampling for efficient completion
~\cite{ju2024brushnet,xie2025turbofill}. Despite strong appearance priors, these
methods define holes in a fixed image coordinate system and do not describe
geometric provenance. Endoscopic completion adds domain cues: temporal learning
removes specularities, DAEVI combines spatio-temporal RGB and depth, and SSIFNet
fuses flow and stereo-temporal information for instrument occlusions
~\cite{daher2023temporal,zhang2024daevi,zou2025ssifnet}. Such methods exploit
neighboring frames or views but still restore defects defined in the original
camera coordinates. Cross-endoscope synthesis instead creates holes through
viewpoint transfer: target-only anatomy is unobserved, while rasterization gaps
or depth conflicts may retain source support. DistriNVS therefore uses
calibrated projection statistics to distinguish true disocclusions from
rasterization gaps and depth conflicts.

\subsection{Geometry-Guided Transport and Completion}

Geometry-guided view synthesis preserves correspondence by transporting source
evidence before generating missing content. SynSin predicts and projects a
feature point cloud, softmax splatting resolves many-to-one mappings with
learned importance, and GeoFill refines depth and pose before reprojection and
inpainting~\cite{wiles2020synsin,niklaus2020softmax,zhao2023geofill}. Recent
feed-forward Gaussian renderers use stronger explicit representations: MVSplat
localizes Gaussians with plane-sweep cost volumes, whereas DepthSplat combines
multi-view correspondence with monocular depth priors
~\cite{chen2024mvsplat,xu2025depthsplat}. They reduce per-scene optimization but
infer a radiance representation from multiple calibrated observations, without
an explicit per-pixel boundary between measured and synthesized appearance.
Generative methods target wider gaps. LeftRefill frames reference-conditioned
synthesis as paired-canvas completion, MultiDiff conditions video diffusion on
depth-warped reference views, and point-guided latent diffusion uses
stereo-derived clouds for laparoscopic view synthesis
~\cite{cao2024leftrefill,muller2024multidiff,tang2025pointguided}. These models can
generate plausible structures in large disocclusions, yet plausibility alone
does not reveal which pixels remain geometrically supported. DistriNVS instead
transports a calibrated source observation with distributional splatting,
derives support, variance, and collision statistics during rasterization, and
uses them to bound synthesis while preserving trusted pixels.

\subsection{Endoscopic Novel-View Synthesis and Uncertainty}

Endoscopic neural rendering must handle non-rigid tissue, view-dependent
illumination, tool occlusions, and sparse camera coverage. EndoNeRF represents
deformation with a neural radiance field, while EndoSurf combines dynamic
modeling with a neural surface~\cite{wang2022endonerf,zha2023endosurf}.
EndoSparse introduces geometric and appearance priors into Gaussian splatting,
and Endo-GSMT uses depth, dense-displacement, and motion priors for dynamic
reconstruction~\cite{li2024endosparse,gou2025endogsmt}. GastroNVS further shows
that large viewpoint changes in real gastroscopy amplify geometric artifacts
and inter-view illumination mismatch~\cite{minai2026gastronvs}. These systems
build persistent scene representations from multiple frames and principally
address reconstruction fidelity within a captured scene. Uncertainty-aware
approaches address sparse evidence from another direction: UC-NeRF conditions a
radiance field on multi-view uncertainty, whereas ExtraGS uses uncertainty to
sample virtual cameras and weight diffusion-refined pseudo-observations for
extrapolation~\cite{guo2025ucnerf,hsieh2026extrags}. Thus, neither uncertainty nor
generative priors are new to endoscopic rendering. The distinction is
operational: DSS propagates a distribution for each input depth through the
calibrated cross-endoscope transform and derives target support, depth variance,
and collision entropy from the same rasterization that transports appearance.
DistriNVS uses these quantities to route completion and protect trusted
measurements during unseen-sequence inference without target-scene optimization.

\section{Method}
\label{sec:method}

\subsection{Problem Formulation and Overview}

Given a synchronized, calibrated cross-endoscope pair, DistriNVS maps a source
RGB image $\I_s$, sparse source depth $\D_s$, its validity mask $\M_s$, camera
intrinsics $\K_s,\K_t$, and the rigid transform
$\T_{s\rightarrow t}\in\mathrm{SE}(3)$ to target RGB
$\widehat\I_t$, target depth $\widehat\D_t$, and a per-pixel risk proxy $\U_t$.
The learned mapping is parameterized by $\theta$. The synchronized target
RGB-D pair and target instrument mask are used only for training losses or
evaluation; none is an inference input.

DistriNVS follows four stages (Fig.~\ref{fig:overview}). First, a depth
reliability encoder converts sparse RGB-D into a source feature map and a
per-pixel depth distribution. Second, distributional surface splatting (DSS)
transports RGB and features while exposing support, depth variance, and
collision entropy in the target view. Third, these renderer statistics control
a mixture of a transport expert and a synthesis expert. Finally, a spatially
consistent trusted mask hard-copies reliable rendered evidence into the output.
A depth-guided backward photometric cycle and the risk objective are active
only during training.

\subsection{Reliability-Aware Source Geometry}

Naive depth completion assigns the same status to measured and inferred
geometry. To retain this distinction, the reliability encoder processes
$[\I_s,\widetilde\D_s,\M_s]$, where $s_D$ is the median valid depth of the
sample and
$\widetilde\D_s=\operatorname{clip}(\D_s/s_D,0,5)$ on valid pixels and zero
elsewhere. During training, $\M_s$ is the post-dropout input mask and
$\M_s^{\rm gt}$ is the pre-dropout supervision mask. The encoder predicts a
source feature $\F_s$, depth mean $\mu$, uncertainty scale $\sigma$,
confidence $c$, and a unit orientation proxy $\n$. For source pixel $q$, let
$m_q=\M_s(q)$. A raw residual $r_q$ parameterizes the depth mean as
\begin{equation}
 \mu_q =
 \begin{cases}
 \D_s(q)\,[1+\eta_D\tanh(r_q)], & m_q=1,\\
 s_D[\operatorname{softplus}(r_q)+\delta_D], & m_q=0,
 \end{cases}
\label{eq:source-distribution}
\end{equation}
Both branches are clipped to the physical interval
$[D_{\min},D_{\max}]$. Here, $\eta_D=0.1$ bounds corrections to observed
measurements and $\delta_D=0.05$ keeps completed depth positive. A second raw
output $r_q^\sigma$ is mapped by a sigmoid to
$[\sigma_D^{\min},\sigma_D^{\max}]$. The confidence range is $[0.5,1]$ for measured pixels
and at most $0.35$ for completed pixels, so inferred geometry can contribute
support but cannot carry measurement-level confidence. The orientation proxy
is normalized with a source-camera optical-axis prior and is learned only
through downstream losses; it is not presented as a metric surface normal.

\subsection{Distributional Surface Splatting}

Given the appearance-feature tensor $\mathbf A_s=[\I_s,\F_s]$, DSS renders
$[\I_t^w,\F_t^w]$ and an expected target depth $\bar D_t$. Unlike a binary
warp, it also returns support $S$, coverage $C$, depth variance $V$, collision
entropy $\mathcal H$, target confidence $\kappa$, and a valid-render mask
$M^{\rm dss}$. Its key distinction from ordinary forward warping is that every
source pixel transports a compact depth distribution rather than a single
point.

\noindent\textbf{Depth hypotheses.}
DSS forms three deterministic hypotheses
$z_{qk}=\operatorname{clip}_{[D_{\min},D_{\max}]}
(\mu_q+\xi_k\sigma_q)$ at $\xi_k\in\{-1,0,1\}$; $a_k$ denotes the fixed
quadrature weight of sample $k$. Each hypothesis is back-projected with
$\K_s$, transformed by $\T_{s\rightarrow t}$, and projected with $\K_t$
under the standard zero-skew pinhole model~\cite{wiles2020synsin,niklaus2020softmax}.
Its target coordinate and target-camera depth are denoted by
$\mathbf x_{qk}=(u_{qk},v_{qk})$ and $z^t_{qk}$, respectively.

\noindent\textbf{Uncertainty-aware visibility.}
Depth uncertainty is converted into a sample-dependent, axis-aligned target
footprint. Let $J^\alpha_{qk}=|\partial\alpha_{qk}/\partial z|$ be the
projected displacement per unit source depth along $\alpha\in\{u,v\}$. Its
scale is
\begin{equation}
 r^\alpha_{qk}=
 \operatorname{clip}_{[0.35,\sigma_{\rm px}^{\max}]}\!\left[
 b+J^\alpha_{qk}\sigma_q+b_n(1-|n_q^z|)\right].
\label{eq:dss-footprint}
\end{equation}
Here, $b$ is the base footprint, $b_n$ expands grazing surfaces, $n_q^z$ is
the optical-axis component of $\n_q$, and $\sigma_{\rm px}^{\max}$ is the
maximum pixel-space scale. For target pixel $p=(u_p,v_p)$, define the
normalized squared offset
$\rho_{qkp}^2=\sum_{\alpha\in\{u,v\}}
[(\alpha_{qk}-\alpha_p)/r^\alpha_{qk}]^2$. The contribution of a projected
hypothesis is
\begin{equation}
 w_{qk\rightarrow p}=c_q a_k\exp\!\left[
 -\frac{\rho_{qkp}^2}{2}
 -\frac{[z^t_{qk}-z_{\min}(p)]_+}{\tau_z}\right].
\label{eq:dss-weight}
\end{equation}
Here, $c_q$ is source confidence, $z_{\min}(p)$ is the detached nearest
arriving depth, $[x]_+=\max(x,0)$, and $\tau_z$ controls soft occlusion. The
minimum and integer assignment are discrete; the remaining quantities are
differentiable. Contributions are evaluated only in the radius-one
neighbourhood of $\lfloor\mathbf x_{qk}\rfloor$; spatial and final weights at
most $10^{-5}$ and $10^{-8}$, respectively, are discarded.

\noindent\textbf{Reliability statistics.}
Let $\Omega_p$ be the set of sample contributions reaching $p$, with
$i=(q_i,k_i)$, source pixel $q_i$, weight $w_i$, and depth $z_i$. Set
$S_p=\sum_{i\in\Omega_p}w_i$ and
$\widetilde S_p=\max(S_p,\epsilon_{\rm dss})$. For any sample attribute $f_i$,
define the DSS aggregation operator
\begin{equation}
 \left\langle f\right\rangle_p^w
 =\sum_{i\in\Omega_p}\frac{w_i}{\widetilde S_p}\,f_i.
\label{eq:dss-moments}
\end{equation}
Applying Eq.~\eqref{eq:dss-moments} to $\mathbf A_s(q_i)$, $z_i$, and
$[z_i-\bar D_t(p)]^2$ gives $\mathbf A_t^w(p)$, $\bar D_t(p)$, and $V_p$,
respectively, where $\mathbf A_t^w=[\I_t^w,\F_t^w]$ and
$\epsilon_{\rm dss}=10^{-8}$. Support is mapped to coverage
$C_p=1-\exp(-S_p)$. Let $\pi_i=w_i/\widetilde S_p$, $N_p=|\Omega_p|$, and
$H_N(\boldsymbol\pi_p)=-\sum_{i\in\Omega_p}\pi_i\log\pi_i/
\log\max(N_p,2)$ denote normalized weight entropy. Because this entropy alone
rises under harmless sub-pixel overlap, DSS modulates it by depth dispersion:
\begin{equation}
 \mathcal H_p=H_N(\boldsymbol\pi_p)
 \left[1-\exp(-V_p/\tau_z^2)\right].
\label{eq:collision-entropy}
\end{equation}
Target confidence $\kappa_p$ is the product of coverage $C_p$, variance
attenuation $\exp[-V_p/(4\tau_z^2)]$, and entropy attenuation
$\exp(-\mathcal H_p)$. The valid-render mask is
$M_p^{\rm dss}=\mathbf 1[S_p\geq S_{\min}]$. Where this mask is zero, DSS
sets $\mathbf A_t^w,\bar D_t,V,\mathcal H,$ and $\kappa$ to zero while
retaining the raw support $S$ and coverage $C$. Thus, $V_p$ and
$\mathcal H_p$ preserve evidence of uncertain depth ordering instead of
reducing reliability to a binary hole mask.

\subsection{Reliability-Guided Completion and Preservation}

Both experts receive $[\I_t^w,\F_t^w,\mathbf h]$, where the six-channel
condition $\mathbf h$ contains coverage, confidence, normalized variance,
collision entropy, normalized target depth, and DSS validity, in that order.
A physics prior $g_p^0$ favours synthesis when confidence falls or geometric
ambiguity rises, and is set to one outside DSS support. We clip this prior by
a margin $\epsilon_g$ and denote its stop-gradient logit by $\ell_p^0$. The
learned router can then make only a bounded correction:
\begin{equation}
 g_p=\operatorname{sigmoid}\!\left[
 \ell_p^0+\eta_g\tanh R_\theta(\mathbf h)_p\right].
\label{eq:router}
\end{equation}
Here, $\epsilon_g$ is the clipping margin, $\eta_g$ bounds the learned
correction, and $R_\theta$ is the residual router. The gate convexly blends
the transport and synthesis RGB-depth predictions: $g_p=0$ selects transport
and $g_p=1$ selects synthesis. We denote the two blended outputs by
$(\I_t^{\rm mix},\D_t^{\rm mix})$ and use $\od$ for element-wise
multiplication below.

The transport expert produces $\I_t^{\rm trans}$ by adding a $\pm0.15$ bounded
RGB residual to $\I_t^w$ and clipping the result to $[0,1]$. The synthesis
expert is a four-resolution U-Net that fuses local
$3\times3$ convolution and real-FFT branches, conditioned at every resolution
by $\mathbf h$. Its five output channels use sigmoid RGB; the depth channel
passes through softplus, receives a $0.05$ positive offset, and is multiplied
by $s_D$; the risk channel uses a sigmoid. These channels produce
$(\I_t^{\rm syn},\D_t^{\rm syn},\U_t^{\rm syn})$. This local--global design
follows Fourier inpainting networks~\cite{suvorov2022lama,chu2023uffc}, but its
weights are trained within DistriNVS rather than imported from either model.
The risk output is treated as an error-scale proxy, not a calibrated failure
probability.

Soft routing alone cannot guarantee preservation. A pixel enters the initial
mask $M^{\rm trust,0}$ only when it has valid DSS support,
$S_p\geq S_{\rm trust}$, $V_p\leq V_{\rm trust}$, and
$\mathcal H_p\leq\mathcal H_{\rm trust}$. We apply a $7\times7$ binary opening
to its untrusted complement and intersect the complement of that result with
$M^{\rm dss}$, yielding the spatially cleaned mask $M^{\rm trust}$. The router
zeros $g$ on $M^{\rm trust,0}$; hard composition then makes the effective gate
zero over the cleaned $M^{\rm trust}$. Let
$\widehat{\mathbf Z}_t=[\widehat\I_t,\widehat\D_t,\U_t]$,
$\mathbf Z_t^w=[\I_t^w,\bar D_t,0]$, and
$\mathbf Z_t^{\rm mix}=[\I_t^{\rm mix},\D_t^{\rm mix},\U_t^{\rm syn}]$,
with binary masks broadcast over channels. The three final outputs share the
single preservation rule
\begin{equation}
 \widehat{\mathbf Z}_t=M^{\rm trust}\od\mathbf Z_t^w+
 (1-M^{\rm trust})\od\mathbf Z_t^{\rm mix}.
\label{eq:hard-compose}
\end{equation}
Thus, trusted output is exactly equal to the accepted DSS warp---not
necessarily ground truth. Zero drift relative to that warp is an architectural
constraint.

Reported inference additionally repairs 8-connected untrusted components that
each occupy at most 2\% of the image. Radius-3 Telea fast-marching
inpainting~\cite{telea2004inpainting} is applied to $\I_t^w$ only on the
recorded repair mask; its values replace the corresponding untrusted output
pixels. Larger components remain learned predictions, trusted pixels are never
changed, and the repair is absent during training.

\subsection{Learning Objectives}

Let $M_{\rm tis}$ exclude target instrument pixels when a tool mask is
available. The untrusted and trusted training regions are
$M_u=(1-M^{\rm trust})\land M_{\rm tis}$ and
$M_v=M^{\rm trust}\land M_{\rm tis}$, respectively. The seam mask $M_b$ is
the morphological boundary of $1-M^{\rm trust}$ generated by a $7\times7$
dilation--erosion window and intersected with $M_{\rm tis}$. Charbonnier reconstruction
terms $\mathcal L_h$, $\mathcal L_v$, and $\mathcal L_b$ act on $M_u$, $M_v$,
and $M_b$, respectively, while $\mathcal L_\nabla$ matches image gradients on
$M_u\cup M_b$. This model-defined training region is distinct from the
raw-warp hole mask used for evaluation.

Available source depth supervises $(\mu,\sigma)$ on the pre-dropout mask
$\M_s^{\rm gt}$ through a masked Laplace negative log-likelihood
$\mathcal L_{d\text{-nll}}$; optional target depth uses
a masked relative-$L_1$ term $\mathcal L_{t\text{-depth}}$. For the
depth-guided cycle, $(\widehat\I_t,\widehat\D_t)$ is rendered back with
$\T_{t\rightarrow s}=\T_{s\rightarrow t}^{-1}$. Detached risk increases the
backward depth scale up to its configured maximum, and confidence is
$1-\U_t$ with a floor of $0.05$. The resulting
$\mathcal L_{\rm cyc}$ is a masked source-view Charbonnier RGB loss; no
separate depth-cycle loss is used.

The router target is one whenever support, variance, or collision entropy
fails its trust threshold and zero otherwise; $\mathcal L_{\rm router}$ is
binary cross-entropy on the pre-sigmoid gate logit. To give risk a statistical
meaning, define the mean RGB absolute error
$e_p=\|\widehat\I_t(p)-\I_t(p)\|_1/3$ and the positive Laplace scale
$s_p=0.01+0.49\U_t(p)$. The risk objective is
\begin{equation}
 \mathcal L_u=\left\langle e_p/s_p+\log s_p\right\rangle_{M_u}.
\label{eq:risk-loss}
\end{equation}
Here, $\I_t$ is the training target and
$\langle\cdot\rangle_M$ denotes the mean over pixels selected by $M$ and, for
channel-valued arguments, over their channels. For a compact presentation,
let $\mathcal L_{\rm app}=\lambda_h\mathcal L_h+\lambda_v\mathcal L_v+
\lambda_b\mathcal L_b+\lambda_\nabla\mathcal L_\nabla$ and
$\mathcal L_{\rm depth}=\mathcal L_{d\text{-nll}}+
\mathcal L_{t\text{-depth}}$. The complete objective is
\begin{equation}
 \mathcal L=\mathcal L_{\rm app}+\lambda_d\mathcal L_{\rm depth}
 +\lambda_c\mathcal L_{\rm cyc}+\lambda_r\mathcal L_{\rm router}
 +\lambda_u\mathcal L_u.
\label{eq:total-loss}
\end{equation}
All $\lambda$ terms are non-negative loss weights; their implemented values
are given in the experimental setup.

The learned inference path ends at Eq.~\eqref{eq:hard-compose}, followed only by
the disclosed deterministic repair in reported runs; target supervision, tool
masks, and the backward cycle are not required. The current formulation assumes
known zero-skew pinhole intrinsics and a rigid inter-endoscope transform, and it uses an
axis-aligned approximation to the projected uncertainty footprint. These
choices define the scope of the method and are tested separately from the
learned completion model in the ablations.

\section{Experiments}
\label{sec:experiments}

\subsection{Datasets and Zero-Shot Protocol}

Training uses the frozen list of 13 iMED Task~2 scene directories supplied in
the project configuration~\cite{imed2026}. Splits are made at scene level, not
at frame level. The training loader uses a synchronized source RGB-depth and
target RGB pair at $640\times512$ resolution. Fifteen percent of valid source
depth samples are removed online, with additional $9\times9$ missing islands,
so that the reliability encoder observes both isolated and contiguous depth
failures. EndoVis/SCARED datasets 8 and 9 are reserved for final zero-shot
evaluation; neither their images nor metrics are used for model selection.

We determine inclusion using a renderer-only pass from the \emph{raw} valid
source depth. A frame is evaluated only if
$|M_{\rm raw}^{\rm dss}|/(H_IW_I)>0.5$; the learned depth completion, target image,
and target tool mask cannot influence this decision. This reproduces the
project's requested strict overlap criterion and prevents a learned module from
selecting favourable frames. The current tables report frame-wise means and
sample standard deviations, both rounded to two decimal places. Paired
sequence-level confidence intervals and permutation tests will be added once
all baseline predictions are frozen.

\todo{Verify the final dataset provenance, license statement, total training
pair count, held-out sequence count, and the immutable list of selected test
frames before submission.}

\subsection{Baselines, Metrics, and Implementation}

Geometry baselines will include hard point projection, nearest-neighbour
small-hole filling, deterministic soft splatting, and softmax splatting. Image
completion baselines will include Telea, LaMa, UFFC, MAT, Stable Diffusion 1.5
Inpainting, and the existing LCM+endoscopy-LoRA system. Directly relevant
endoscopic baselines (DAEVI and SSIFNet) will be evaluated when code and an
equivalent cross-view protocol are available; otherwise their non-comparable
reported values will not be inserted into the main table. Every completion
method will use the same raw mask, spatial resolution, selected frames, and
known-pixel hard composition.

The existing evaluator reports full-, raw-hole-, raw-visible-, and 7-pixel
seam-band PSNR/SSIM, full-image and hole-region LPIPS, known-region drift, risk
expected calibration error (ECE), area under the risk--coverage curve (AURC),
and a projection-defect taxonomy. The final study will additionally report DISTS, depth AbsRel/RMSE,
disparity endpoint error, normal error, latency, throughput, peak memory, and
parameter count. Results will be stratified by raw overlap, hole ratio, depth
validity, instrument presence, and sequence.

The reference configuration uses sigma offsets $(-1,0,1)$ with quadrature
weights $(0.2741,0.4518,0.2741)$, a one-pixel splat radius,
$\tau_z=2$ mm, $S_{\min}=0.03$, $S_{\rm trust}=0.35$,
$V_{\rm trust}=9$ mm$^2$, $\mathcal H_{\rm trust}=1.0$, a 7-pixel seam band,
and base width 32. The reliability encoder combines a full-resolution residual
branch with a quarter-resolution context branch; when no valid source depth is
available, it uses $s_D=100$ mm. Depth and uncertainty are bounded to $[1,500]$ mm and
$[0.25,20]$ mm, respectively; $b=b_n=0.75$ px,
$\sigma_{\rm px}^{\max}=2.5$ px, and the Charbonnier constant is $10^{-3}$.
The router uses $\bar V=V/(V+\tau_z^2)$ and
$\widetilde D=\operatorname{clip}(\bar D_t/s_D,0,5)/5$ as its normalized
variance and depth channels. Its prior is
$g^0=0.65(1-\kappa)+0.20\bar V+0.15\mathcal H$, with
$\epsilon_g=10^{-4}$ and $\eta_g=0.25$.
Because $\mathcal H\in[0,1]$, this reference threshold is
neutral in the hard trusted mask; collision entropy still affects
$\kappa$ and the soft synthesis gate. Reported inference uses the disclosed
2\%-area, radius-3 Telea repair for small untrusted components. Training is
scheduled for 10,000 optimizer steps with AdamW, learning rate $2\times10^{-4}$,
weight decay $10^{-4}$, and three-GPU distributed training. The per-GPU batch
size is 1 with four-step gradient accumulation, giving an effective batch size
of 12. We use mixed precision, a maximum gradient norm of 1, and seed 6666. The
implemented loss weights
$(\lambda_h,\lambda_v,\lambda_b,\lambda_\nabla,\lambda_d,\lambda_c,
\lambda_r,\lambda_u)$ are $(3,1,2,0.2,0.2,0.2,0.1,0.1)$.

\todo{Record the final checkpoint, model-selection rule, GPU model, training
time, inference warm-up, precision, exact parameter/FLOP counts, and all
baseline versions and checkpoints.}

\subsection{Main Comparison}

We first compared all methods on the same 1,392-frame iMED validation set. To
avoid method-dependent masks, every prediction was re-evaluated against the
same target image and raw-warp hole mask. The mask covered 25.84\% of an image
on average. We additionally re-evaluated four scene-optimized Gaussian
renderers from the comparison benchmark; each retained method produced a
prediction for every selected frame. These Gaussian methods use the
dataset-provided ground-truth depth and are therefore reference reconstruction
baselines rather than matched feed-forward completion models. PR-ENDO,
EndoGS, and EndoGaussian were omitted because their available predictions cover
only 1/7, 5/7, and 6/7 validation scenes, respectively.
Table~\ref{tab:validation-comparison} reports the frame-wise mean and sample
standard deviation under this common protocol.
The deterministic small-component repair affected
$1.17\pm0.94$\% of image pixels per validation frame on average; the 2\%
criterion applies to each connected component, so their union may be larger.

\begin{table*}[t]
\centering
\caption{Unified comparison on the 1,392-frame iMED validation set. All methods
use identical targets and raw-warp hole masks. Values are frame-wise
mean~$\pm$~SD; best means are bold. Gaussian baselines are scene-optimized
using dataset ground-truth depth and have complete 7/7-scene coverage. The
DistriNVS row includes the disclosed small-component Telea repair.}
\label{tab:validation-comparison}
\setlength{\tabcolsep}{3.5pt}
\begin{tabular}{lcccccc}
\toprule
Method & PSNR $\uparrow$ & SSIM $\uparrow$ & LPIPS $\downarrow$ &
H-PSNR $\uparrow$ & H-SSIM $\uparrow$ & H-LPIPS $\downarrow$ \\
\midrule
\multicolumn{7}{l}{\textit{Scene-optimized Gaussian reconstruction}} \\
Deform3DGS~\cite{yang2024deform3dgs} & $15.55\pm1.78$ & $0.55\pm0.04$ &
$0.63\pm0.06$ & $14.11\pm1.37$ & $0.47\pm0.10$ & $0.69\pm0.06$ \\
Endo-4DGS~\cite{huang2024endo4dgs} & $16.44\pm1.71$ & $0.58\pm0.04$ &
$0.61\pm0.06$ & $15.61\pm1.47$ & $0.57\pm0.13$ & $0.66\pm0.06$ \\
Free-SurGS~\cite{guo2024freesurgs} & $16.83\pm2.11$ & $0.67\pm0.06$ &
$0.63\pm0.04$ & $17.05\pm2.38$ & $0.72\pm0.04$ & $0.61\pm0.04$ \\
SurgicalGS~\cite{chen2025surgicalgs} & $17.41\pm1.99$ & $0.60\pm0.06$ &
$0.60\pm0.05$ & $16.98\pm2.55$ & $0.61\pm0.14$ & $0.65\pm0.04$ \\
\addlinespace
\multicolumn{7}{l}{\textit{Image completion}} \\
SD1.5 Inpainting & $16.84\pm2.73$ & $0.57\pm0.05$ & $0.52\pm0.06$ &
$14.56\pm3.14$ & $0.50\pm0.10$ & $0.64\pm0.07$ \\
LaMa & $18.99\pm1.58$ & $0.64\pm0.07$ & $\mathbf{0.47\pm0.03}$ &
$20.41\pm2.06$ & $0.74\pm0.05$ & $\mathbf{0.48\pm0.04}$ \\
LCM + endoscopy LoRA & $19.02\pm1.64$ & $0.63\pm0.09$ & $0.49\pm0.03$ &
$19.54\pm2.28$ & $0.74\pm0.09$ & $0.49\pm0.03$ \\
\midrule
DistriNVS (ours) & $\mathbf{19.96\pm1.79}$ & $\mathbf{0.69\pm0.08}$ &
$0.55\pm0.03$ & $\mathbf{23.98\pm3.28}$ & $\mathbf{0.82\pm0.06}$ &
$0.52\pm0.03$ \\
\bottomrule
\end{tabular}
\end{table*}

DistriNVS achieved the best distortion- and structure-based result, improving
full-image PSNR by 0.94~dB and SSIM by 0.05 over the strongest corresponding
image-completion baselines. When the Gaussian reconstruction baselines are included, the SSIM
margin narrows to 0.02 over Free-SurGS, whereas the PSNR margin remains
0.94~dB over the LCM baseline. The advantage was larger inside the missing region: hole PSNR
increased by 3.57~dB over LaMa, while hole SSIM increased by 0.08 over the
supervised LoRA baseline. However, LaMa obtained lower full-image LPIPS
($0.47\pm0.03$ versus $0.55\pm0.03$) and lower hole LPIPS
($0.48\pm0.04$ versus $0.52\pm0.03$). The results therefore show stronger
pixel and structural fidelity for DistriNVS, but not a uniform advantage in
learned perceptual distance.

We separately evaluated DistriNVS on held-out EndoVis Datasets~8 and 9 using
the strict $\mathrm{raw\ overlap}>0.5$ inclusion rule. Among 124 candidate
frames, 78 met the rule. Table~\ref{tab:zeroshot-breakdown} reports the two
datasets separately and jointly.

\begin{table*}[t]
\centering
\caption{DistriNVS zero-shot performance on held-out EndoVis Datasets~8 and 9.
Values are frame-wise mean~$\pm$~SD over frames satisfying
$\mathrm{raw\ overlap}>0.5$ and include the disclosed small-component repair.}
\label{tab:zeroshot-breakdown}
\setlength{\tabcolsep}{3.2pt}
\begin{tabular}{lrrrrrrr}
\toprule
Dataset & Frames & PSNR $\uparrow$ & SSIM $\uparrow$ & LPIPS $\downarrow$ &
H-PSNR $\uparrow$ & H-SSIM $\uparrow$ & H-LPIPS $\downarrow$ \\
\midrule
EndoVis 8 & 48 & $\mathbf{24.85\pm2.00}$ & $\mathbf{0.72\pm0.12}$ &
$\mathbf{0.21\pm0.04}$ & $\mathbf{23.71\pm2.74}$ &
$\mathbf{0.70\pm0.09}$ & $\mathbf{0.23\pm0.06}$ \\
EndoVis 9 & 30 & $22.10\pm2.29$ & $0.63\pm0.09$ & $0.27\pm0.06$ &
$18.80\pm3.32$ & $0.45\pm0.08$ & $0.35\pm0.08$ \\
\midrule
Combined & 78 & $23.79\pm2.49$ & $0.68\pm0.12$ & $0.23\pm0.06$ &
$21.82\pm3.81$ & $0.61\pm0.15$ & $0.28\pm0.09$ \\
\bottomrule
\end{tabular}
\end{table*}

Across the 78 accepted frames, DistriNVS obtained
$23.79\pm2.49$~dB PSNR, $0.68\pm0.12$ SSIM,
$0.23\pm0.06$ LPIPS, $21.82\pm3.81$~dB hole PSNR,
$0.61\pm0.15$ hole SSIM, and $0.28\pm0.09$ hole LPIPS. The deterministic
repair affected $3.75\pm2.11$\% of pixels per accepted frame. Known-region
drift was zero under the model-defined trusted mask. Baseline inference on
Datasets~8 and 9 is not yet complete; consequently,
Table~\ref{tab:zeroshot-breakdown} does not support a cross-method zero-shot
ranking. Paired confidence intervals and significance tests will be added after
the final prediction set is frozen.

\subsection{Ablations and Geometric Diagnostics}

The core ablation in Table~\ref{tab:ablation} follows the implementation flags.
The hard-point variant disables depth completion, distributional samples,
normal-dependent footprint, physics routing, and the backward cycle. The soft
splat variant retains a deterministic Gaussian footprint. Subsequent rows add
DSS, the physics router, hard composition, and the RGB-D-defined backward
cycle. This ordering tests whether an improvement originates from ordinary
anti-aliased splatting or from the proposed distributional visibility model.

\begin{table}[t]
\centering
\caption{Component ablation. $K$ is the number of depth sigma points.}
\label{tab:ablation}
\setlength{\tabcolsep}{3.2pt}
\begin{tabular}{lccccc}
\toprule
Variant & $K$ & DSS & Router & Hard & PSNR \\
\midrule
Hard point & 1 & -- & -- & \checkmark & \tbd\\
Soft splat & 1 & -- & -- & \checkmark & \tbd\\
DSS only & 3 & \checkmark & -- & \checkmark & \tbd\\
DSS + router & 3 & \checkmark & \checkmark & \checkmark & \tbd\\
No hard composition & 3 & \checkmark & \checkmark & -- & \tbd\\
No backward cycle & 3 & \checkmark & \checkmark & \checkmark & \tbd\\
Full DistriNVS & 3 & \checkmark & \checkmark & \checkmark & \tbd\\
\bottomrule
\end{tabular}
\end{table}

\todo{Add columns for seam error, visible-region PSNR, known drift, depth
error, AURC, and runtime. Run focused DSS ablations for 1/3/5 sigma
points, isotropic versus anisotropic footprints, normal removal, hard versus
soft visibility, and individual $S/V/\mathcal H$ router inputs.}

\todo{Report the paired with/without small-component Telea repair ablation;
keep the repair disabled in all architecture-only
ablations.}

To determine whether DSS solves the intended geometric problem before learned
completion is considered, we will report raw coverage, visible-region RGB
error, target-depth error, and occlusion-boundary ordering for hard, softmax,
and distributional splatting. The diagnostic label map partitions target pixels
into reliable/unclassified support, local sampling gaps, depth-completion-
recoverable gaps, unsupported/disoccluded regions, and surface-collision risk.
These labels are diagnostics rather than ground truth because a single source
view cannot perfectly separate genuine disocclusion from catastrophic depth
failure.

\todo{Insert the DSS-only table and plots showing monotonic relationships
between support/variance/collision entropy and actual reprojection error.}

\subsection{Qualitative, Failure, and Efficiency Analysis}

\begin{figure*}[t]
\centering
\fbox{\parbox[c][35mm][c]{0.96\textwidth}{\centering
\todo{Replace this placeholder with synchronized columns: source RGB, raw hard
warp, raw validity, LaMa, LCM, DistriNVS, target RGB, support, variance,
collision entropy, and risk. Include at least one success and one failure for
each overlap bin.}}}
\caption{Planned qualitative comparison and failure visualization. Crops and
color scales must be identical across methods; generated and source-backed
regions will be explicitly marked.}
\label{fig:qualitative}
\end{figure*}

Efficiency is central to the ICRA framing. We will measure end-to-end latency
including depth reliability estimation, DSS, both experts, and hard composition;
DSS-only timing will also be reported to isolate the renderer cost. Failure
analysis will cover specular depth loss, calibration perturbation, thin
instruments, large disocclusion, weak texture, and dynamic tissue deformation.

\todo{Insert the speed/memory/parameter table, calibration-perturbation curves,
and blinded expert assessment if available. Do not claim real-time performance
until measured on the stated hardware and resolution.}

\section{Discussion}
\label{sec:discussion}

Under a common validation protocol, DistriNVS outperformed the evaluated
Gaussian reconstruction, generic inpainting, and supervised LoRA baselines in
both full-image and hole-region distortion and structural fidelity. The
Gaussian rows provide a reference rather than a matched training comparison,
because they are optimized per scene with dataset ground-truth depth. The
larger improvement in hole PSNR (3.57~dB) than in full-image PSNR (0.94~dB) is
consistent with the intended separation between evidence
transport and unsupported-region synthesis. Together with zero measured
drift under the trusted mask, these observations support the operational value
of constraining where synthesis is permitted. They do not yet isolate the
individual contributions of distributional splatting, routing, and hard
composition; the planned component ablations remain necessary for that causal
claim. Moreover, DistriNVS's higher LPIPS than LaMa shows that its gains in
pixel and structural agreement do not translate uniformly to the AlexNet
perceptual feature space; qualitative analysis is needed to determine whether
this difference reflects texture realism, colour statistics, or metric bias.

The method does not remove the fundamental ambiguity of unseen tissue. A low
risk score indicates model confidence under the training distribution; it does
not certify anatomical truth. Moreover, the current formulation assumes known
intrinsics and a rigid inter-camera transform for each pair, processes frames
independently, and uses an axis-aligned rather than full-covariance footprint.
The zero-shot evaluation also revealed a gap between the two held-out datasets:
full-image PSNR decreased from $24.85\pm2.00$~dB on EndoVis~8 to
$22.10\pm2.29$~dB on EndoVis~9, while hole SSIM decreased from
$0.70\pm0.09$ to $0.45\pm0.08$. This difference limits claims of uniform
cross-dataset generalization, and incomplete baseline runs currently prevent a
fair cross-method ranking on Datasets~8 and 9.
Large calibration errors, non-rigid motion between nominally synchronized
frames, or out-of-distribution anatomy may therefore invalidate the transport
statistics. Generated output should remain visibly marked and should not be
used as a substitute for a live physical endoscopic view.

\todo{Rewrite the first discussion paragraph from the final ablation and risk
curves. Add observed failure modes, not generic limitations, and compare them
directly with LaMa/UFFC, DAEVI/SSIFNet, UC-NeRF, and ExtraGS.}

\section{Conclusion}
\label{sec:conclusion}

This paper introduces a reliability-aware formulation of cross-endoscope
novel-view synthesis and implements it through distributional surface
splatting, visibility-constrained transport/synthesis routing, trusted-pixel
hard composition, and joint appearance--geometry risk modeling. The design
turns projection support, depth variance, and collision entropy into operational
constraints on what the network may synthesize. \todo{After the frozen
evaluation, replace the following validation evidence with the frozen test-set
comparison if available.} On a common 1,392-frame validation protocol,
DistriNVS exceeded the strongest per-metric baselines by 0.94~dB in
full-image PSNR and 3.57~dB in hole PSNR, while attaining a hole SSIM of
$0.82\pm0.06$.
These results support visibility-constrained completion under the evaluated
setting, although LaMa retained the lowest LPIPS. The difference between
EndoVis Datasets~8 and 9 and unfinished zero-shot baseline evaluation define
the current boundary of the evidence.

\bibliographystyle{IEEEtran}
\bibliography{references}

\end{document}
